\documentclass[10pt,a4paper]{amsart}
\usepackage{amsmath,amssymb,amsthm,graphicx}
% Retained minimal preamble for the Stacks Project chapter flat.tex (More on Flatness).
% Verbatim excerpt of the chapter preamble: Xy-pic setup (native preamble lines 14-18) and the native
% theorem declarations with the mathematical macros (native preamble lines 154-217). Copied mechanically;
% nothing added, removed or redefined. The wrapper class and the amsmath/amssymb/amsthm/graphicx base
% packages are supplied by the assembler, so the class/externaldocument part of the native preamble is omitted.
% For commutative diagrams we use Xy-pic
\usepackage[all]{xy}

% We use 2cell for 2-commutative diagrams.
\xyoption{2cell}
% Theorem environments.
%
\theoremstyle{plain}
\newtheorem{theorem}[subsection]{Theorem}
\newtheorem{proposition}[subsection]{Proposition}
\newtheorem{lemma}[subsection]{Lemma}

\theoremstyle{definition}
\newtheorem{definition}[subsection]{Definition}
\newtheorem{example}[subsection]{Example}
\newtheorem{exercise}[subsection]{Exercise}
\newtheorem{situation}[subsection]{Situation}

\theoremstyle{remark}
\newtheorem{remark}[subsection]{Remark}
\newtheorem{remarks}[subsection]{Remarks}

\numberwithin{equation}{subsection}

% Macros
%
\def\lim{\mathop{\mathrm{lim}}\nolimits}
\def\colim{\mathop{\mathrm{colim}}\nolimits}
\def\Spec{\mathop{\mathrm{Spec}}}
\def\Hom{\mathop{\mathrm{Hom}}\nolimits}
\def\Ext{\mathop{\mathrm{Ext}}\nolimits}
\def\SheafHom{\mathop{\mathcal{H}\!\mathit{om}}\nolimits}
\def\SheafExt{\mathop{\mathcal{E}\!\mathit{xt}}\nolimits}
\def\Sch{\mathit{Sch}}
\def\Mor{\mathop{\mathrm{Mor}}\nolimits}
\def\Ob{\mathop{\mathrm{Ob}}\nolimits}
\def\Sh{\mathop{\mathit{Sh}}\nolimits}
\def\NL{\mathop{N\!L}\nolimits}
\def\CH{\mathop{\mathrm{CH}}\nolimits}
\def\proetale{{pro\text{-}\acute{e}tale}}
\def\etale{{\acute{e}tale}}
\def\QCoh{\mathit{QCoh}}
\def\Ker{\mathop{\mathrm{Ker}}}
\def\Im{\mathop{\mathrm{Im}}}
\def\Coker{\mathop{\mathrm{Coker}}}
\def\Coim{\mathop{\mathrm{Coim}}}

% Boxtimes
%
\DeclareMathSymbol{\boxtimes}{\mathbin}{AMSa}{"02}

%
% Macros for moduli stacks/spaces
%
\def\QCohstack{\mathcal{QC}\!\mathit{oh}}
\def\Cohstack{\mathcal{C}\!\mathit{oh}}
\def\Spacesstack{\mathcal{S}\!\mathit{paces}}
\def\Quotfunctor{\mathrm{Quot}}
\def\Hilbfunctor{\mathrm{Hilb}}
\def\Curvesstack{\mathcal{C}\!\mathit{urves}}
\def\Polarizedstack{\mathcal{P}\!\mathit{olarized}}
\def\Complexesstack{\mathcal{C}\!\mathit{omplexes}}
% \Pic is the operator that assigns to X its picard group, usage \Pic(X)
% \Picardstack_{X/B} denotes the Picard stack of X over B
% \Picardfunctor_{X/B} denotes the Picard functor of X over B
\def\Pic{\mathop{\mathrm{Pic}}\nolimits}
\def\Picardstack{\mathcal{P}\!\mathit{ic}}
\def\Picardfunctor{\mathrm{Pic}}
\def\Deformationcategory{\mathcal{D}\!\mathit{ef}}

% --- appended 2026-10-12 (Chief): the project own definitions of the reference/slogan/history
% display environments, copied mechanically from the shared Stacks preamble lines 8-12
% (\usepackage{verbatim} plus the three \newenvironment lines, which the chapter preamble keeps
% verbatim from the project preamble). Nothing added, removed or redefined: these environments
% carry no mathematical content and the project itself typesets them as comments. Without these
% three lines a counted statement whose own text uses them fails with Environment slogan undefined.
\usepackage{verbatim}
\newenvironment{reference}{\comment}{\endcomment}
\newenvironment{slogan}{\comment}{\endcomment}
\newenvironment{history}{\comment}{\endcomment}
\usepackage{cleveref}
\crefname{originallemma}{original lemma}{original lemmas}
\Crefname{originallemma}{Original lemma}{Original lemmas}
\makeatletter
\newlabel{more-algebra-lemma-limit-preserving-flat-at-primes}{{05LT}{source}{}{sourcepointer-1}{}}
\newlabel{more-algebra-lemma-limit-preserving-flat-at-primes@cref}{{[originallemma][0][]05LT}{[1][0][]source}{}{}{}}
\makeatother
\crefname{originalremark}{original remark}{original remarks}
\Crefname{originalremark}{Original remark}{Original remarks}
\makeatletter
\newlabel{remark-flattening-local-scheme-theoretic}{{05PH}{source}{}{sourcepointer-2}{}}
\newlabel{remark-flattening-local-scheme-theoretic@cref}{{[originalremark][0][]05PH}{[1][0][]source}{}{}{}}
\makeatother
\crefname{originaltheorem}{original theorem}{original theorems}
\Crefname{originaltheorem}{Original theorem}{Original theorems}
\makeatletter
\newlabel{theorem-flattening-local}{{05PG}{source}{}{sourcepointer-3}{}}
\newlabel{theorem-flattening-local@cref}{{[originaltheorem][0][]05PG}{[1][0][]source}{}{}{}}
\makeatother
\title{Flattening over the points of a given subset of the closed fibre is cut out by a closed subscheme}
\author{The Stacks Project Authors}
\date{stacks-project:flat.tex:local-snapshot-20261011; GNU FDL 1.2 or later}
\begin{document}
\maketitle
\noindent\textit{Editorial note.} All delivered mathematical text is the author's own, copied from the pinned Stacks Project chapter file flat.tex; no statement or proof step was rewritten, reordered or removed.\par\medskip
\noindent\textit{Editorial note.} Preamble provenance: the chapter's own retained minimal preamble (sources/stacks-flat/retained-preamble.tex, lines 1-85) is a verbatim excerpt of the native Xy-pic setup and of the native theorem declarations and mathematical macros, followed by the project's own reference/slogan/history display-environment definitions copied mechanically from the shared preamble; the wrapper class and the amsmath/amssymb/amsthm/graphicx base packages are supplied by the assembler, and extra\_packages is left empty because everything needed is already in the retained preamble. No macro was invented or redefined.\par\medskip
\noindent\textit{Editorial note.} Same-chapter pointers are located at their \textbackslash{}label lines inside the pinned chapter file; cross-chapter pointers are located in the project's own tag list staged in the snapshot (sources/stacks-chow/tags.txt), at the verbatim lines carrying their canonical tags.\par\medskip
\noindent\textit{Allowed standard prerequisites (editorial).}\begin{itemize}
\item The scheme theoretic version of the local flattening statement. Same chapter flat.tex, label remark-flattening-local-scheme-theoretic, invoked at proof line 6950. Provides, for each point of the closed fibre, a closed subscheme of the base whose effect is exactly to force flatness at that point, so that the intersection over the given subset does the job.
\item Limit preserving behaviour of the flatness functor at primes, and the finite presentation argument from the local flattening theorem. More on Algebra, label more-algebra-lemma-limit-preserving-flat-at-primes (proof line 6961); Same chapter flat.tex, the final part of the proof of theorem-flattening-local referred to at proof lines 6963-6965. The first shows the functor is limit preserving when the sources and the module are of finite presentation, and the second supplies the argument giving finite presentation of the closed subscheme.
\end{itemize}\medskip
\ifdefined\hypertarget\hypertarget{sourcepointer-1}{}\fi
\noindent\textit{Source reference.} Original lemma 05LT: sources/stacks-chow/tags.txt, source lines 7247--7247. Navigation pointer for \textbackslash{}ref\{more-algebra-lemma-limit-preserving-flat-at-primes\} used inside the counted span (line 6961); the argument that consumes it is retained in the span.\par\medskip
\ifdefined\hypertarget\hypertarget{sourcepointer-2}{}\fi
\noindent\textit{Source reference.} Original remark 05PH: sources/stacks-flat/pinned-original-source.tex, source lines 6905--6905. Navigation pointer for \textbackslash{}ref\{remark-flattening-local-scheme-theoretic\} used inside the counted span (line 6950); the argument that consumes it is retained in the span.\par\medskip
\ifdefined\hypertarget\hypertarget{sourcepointer-3}{}\fi
\noindent\textit{Source reference.} Original theorem 05PG: sources/stacks-flat/pinned-original-source.tex, source lines 6790--6790. Navigation pointer for \textbackslash{}ref\{theorem-flattening-local\} used inside the counted span (line 6965); the argument that consumes it is retained in the span.\par\medskip
\begin{lemma}
\label{lemma-freebie}
Let $S$ be the spectrum of a henselian local ring with closed point $s$.
Let $X \to S$ be a morphism of schemes which is locally of finite type.
Let $\mathcal{F}$ be a finite type quasi-coherent $\mathcal{O}_X$-module.
Let $E \subset X_s$ be a subset. There exists a closed subscheme
$Z \subset S$ with the following property: for any morphism of pointed
schemes $(T, t) \to (S, s)$ the following are equivalent
\begin{enumerate}
\item $\mathcal{F}_T$ is flat over $T$ at all points of the fibre
$X_t$ which map to a point of $E \subset X_s$, and
\item $\Spec(\mathcal{O}_{T, t}) \to S$ factors through $Z$.
\end{enumerate}
Moreover, if $X \to S$ is locally of finite presentation,
$\mathcal{F}$ is of finite presentation, and $E \subset X_s$ is
closed and quasi-compact, then $Z \to S$ is of finite presentation.
\end{lemma}

\begin{proof}
For $x \in X_s$ denote $Z_x \subset S$ the closed subscheme we found in
Remark \ref{remark-flattening-local-scheme-theoretic}.
Then it is clear that $Z = \bigcap_{x \in E} Z_x$ works!

\medskip\noindent
To prove the final statement assume $X$ locally of finite presentation,
$\mathcal{F}$ of finite presentation and $Z$ closed and quasi-compact.
First, choose finitely many affine opens $W_j \subset X$ such that
$E \subset \bigcup W_j$. It clearly suffices to prove the
result for each morphism $W_j \to S$ with sheaf $\mathcal{F}|_{X_j}$
and closed subset $E \cap W_j$. Hence we may assume $X$ is affine.
In this case,
More on Algebra, Lemma \ref{more-algebra-lemma-limit-preserving-flat-at-primes}
shows that the functor defined by (1) is ``limit preserving''.
Hence we can show that $Z \to S$ is of finite presentation exactly
as in the last part of the proof of
Theorem \ref{theorem-flattening-local}.
\end{proof}

\end{document}
